%% ---- 3. Formulação do Problema ---------------------------------------------
\section{Formulação do Problema}
\label{sec:formulacao}

\paragraph{Margem.}
O recomendador atribui um escore $f_i(s)$ a cada item $i$ dado o histórico
$s=(s_1,\dots,s_n)$. Fixamos um item $v$ e medimos quão dentro da lista ele
está pela margem
\begin{equation}
  \delta(v \mid s') \;=\; f_v(s') \;-\; f_{(K)}(s'),
  \label{eq:margem}
\end{equation}
em que $f_{(K)}(s')$ é o $K$-ésimo maior escore entre os demais itens de um
conjunto de candidatos fixo $\mathcal{C}$: o catálogo sem os itens que o
usuário já consumiu. O conjunto é o mesmo para todas as intervenções, para que
trocar itens do histórico não altere quem concorre com $v$. Assim,
$v$ está no top-$K$ de $s'$ se e somente se $\delta(v \mid s') \ge 0$. Usamos
$K=10$.

\paragraph{Explicação e intervenções.}
Uma explicação é um conjunto de posições $X \subseteq \{1,\dots,n\}$,
codificado pela máscara $z \in \{0,1\}^n$ com $z_j = 1 \iff j \in X$. Dado um
histórico de fundo $r$ de mesmo comprimento, definimos
\begin{equation}
  s[X \leftarrow r]_j =
  \begin{cases} r_j, & j \in X \\ s_j, & j \notin X \end{cases}
  \label{eq:int-nec}
\end{equation}
\begin{equation}
  s[\bar X \leftarrow r]_j =
  \begin{cases} s_j, & j \in X \\ r_j, & j \notin X \end{cases}
  \label{eq:int-suf}
\end{equation}
A primeira apaga a explicação e mantém o resto; a segunda mantém a explicação
e troca todo o resto. As duas são substituições: o comprimento e as posições
não mudam.

\paragraph{Necessidade e suficiência.}
Com o fundo $R$ sorteado de uma distribuição $\mathcal{R}$, as probabilidades
de necessidade e de suficiência~\cite{pearl2009causality,tan2021counter} são
\begin{equation}
  \PN(X) \;=\; \Pr_{R \sim \mathcal{R}}\!\big[\,\delta(v \mid s[X \leftarrow R]) < 0\,\big]
  \label{eq:pn}
\end{equation}
\begin{equation}
  \PS(X) \;=\; \Pr_{R \sim \mathcal{R}}\!\big[\,\delta(v \mid s[\bar X \leftarrow R]) \ge 0\,\big]
  \label{eq:ps}
\end{equation}

\paragraph{Problema.}
Buscamos as máscaras não dominadas de
\begin{equation}
  \min_{z \in \{0,1\}^n} \;\; \big(\, N(z),\; S(z),\; C(z) \,\big)
  \label{eq:problema}
\end{equation}
\begin{equation}
  N(z) = \E_{R}\big[\,\delta(v \mid s[X \leftarrow R])\,\big]
  \label{eq:N}
\end{equation}
\begin{equation}
  S(z) = -\,\E_{R}\big[\,\delta(v \mid s[\bar X \leftarrow R])\,\big]
  \label{eq:S}
\end{equation}
\begin{equation}
  C(z) = |X|
  \label{eq:C}
\end{equation}
$N$ é a margem esperada de $v$ depois de apagar $X$ (quanto menor, mais
necessário); $S$ é a margem esperada mantendo só $X$, com sinal trocado (quanto
menor, mais suficiente); $C$ é o tamanho. Otimizamos margens médias, e não
$\PN$ e $\PS$, porque estas são quase sempre $0$ ou $1$ e não dão pressão de
seleção à busca; $\PN$ e $\PS$ entram só na avaliação.

\paragraph{Por que é difícil.}
O espaço tem $2^n$ máscaras, com $n$ até 128. O modelo é uma caixa-preta e
cada avaliação de uma máscara custa várias consultas. Os objetivos não são
comensuráveis: o histórico inteiro é ótimo em $N$ e $S$ e pior em $C$. Se $N$
e $S$ também conflitam entre si é uma pergunta empírica, respondida na
Seção~\ref{sec:conflito}.
